\chapter{What Failed in the Particle Proposals, and What Conditioning Repairs}
\label{ch:bf-dpf-observation-parent-failures}
\chaptermark{Particle Proposal Failures}

An observation can improve a proposal in two places: it can move a child
toward likely states, and it can change which parents receive children.
Chapter~\ref{ch:bf-ukf-covariance-proposal} develops the proposal densities
and their importance correction. We now ask what those changes accomplish
for the likelihood and its parameter derivative. The experiments give a
clear reason to retain weighted parents for weight balance, but they also
show why high ESS and a correctly differentiated finite calculation are
insufficient tests of model-score accuracy.

\section{Conditioning the child and allocating the parents}
Let the incoming equally weighted cloud be $a_1,\ldots,a_N$. Its predictive
mixture is $\widehat p^-(x)=N^{-1}\sum_j f_\theta(x\mid a_j)$. Define each
parent's predictive observation density and conditional child density by
\begin{equation}
 L_j(y)=\int f_\theta(x\mid a_j)g_\theta(y\mid x)\,\mathrm dx,
 \qquad
 q_j(x\mid y)=\frac{f_\theta(x\mid a_j)g_\theta(y\mid x)}{L_j(y)}.
\end{equation}
These expressions assume positive finite $L_j$ and evaluable densities.
Multiplying the predictive mixture by the likelihood and normalizing gives
\begin{equation}
 \widehat p(x\mid y)=\sum_j\alpha_j(y)q_j(x\mid y),
 \qquad \alpha_j(y)=\frac{L_j(y)}{\sum_k L_k(y)},
 \qquad \widehat Z=\frac1N\sum_jL_j(y).
\label{eq:parent-finite-mixture-conditioning}
\end{equation}
Thus the observation affects both the conditional state and the probability
of its ancestor. Uniform parent selection uses the same conditional children
but assigns mass $1/N$ to each component. It remains a valid importance
proposal when its actual mixture density is used in the denominator.
An observation-informed UKF guide or LEDH movement alone does not perform
this reallocation of mass between ancestors.

For a Gaussian transition $f(x\mid a_j)=\mathcal N(x;m_j,Q_j)$ and linear
observation $y=Hx+\varepsilon$, $\varepsilon\sim\mathcal N(0,R)$,
\begin{align}
 S_j&=HQ_jH^\top+R,&L_j(y)&=\mathcal N(y;Hm_j,S_j),\\
 K_j&=Q_jH^\top S_j^{-1},&
 q_j&=\mathcal N\!\left(m_j+K_j(y-Hm_j),\,Q_j-K_jS_jK_j^\top\right).
\end{align}
The covariance is conditional on parent $j$. Between-parent uncertainty
belongs to the mixture, rather than being inserted into every $Q_j$.
These identities explain the SIR conditional calculation in its stated
Gaussian observation setting. KSC additionally conditions the observation
mixture component; range--bearing instead uses a local Gaussian approximation
with wrapped bearing residuals. Exactness must be asserted separately for
each target.

\section{Likelihood, score and ESS in the completed comparisons}
The completed comparisons answer two different questions. Conditioning a child
on its parent's transition and the current observation changes where that child
is drawn. Weighting the parents changes how many children are allocated to each
part of the incoming cloud. The original three-branch proposal is the comparator
for the first change; conditional proposals with uniform parents isolate the
second. The conditional-mixture arms retain their defensive branches. Thus
uniform parents do not mean an observation-blind proposal: their children
already use the observation.

The five-model comparison uses $N=1008$, eight paired random streams and the
saved FP64 GPU/XLA diagnostic runs. Each model has one fixed dataset and
parameter configuration; $T=50$ except for range--bearing, where $T=20$.
The SIR row is the nominal $T=50$ case, while the larger SIR campaign varies
parameters and uses $T=10,20,50$. These settings describe the completed
experiment. They are not independently tuned defaults for every model or
evidence for production-precision accuracy.

For a scalar quantity $a$ with reference $a_\star$, the tables report
\begin{equation}
 \overline a=\frac18\sum_{r=1}^8\widehat a_r,
 \qquad
 \operatorname{RMSE}(a)=
 \left\{\frac18\sum_{r=1}^8(\widehat a_r-a_\star)^2\right\}^{1/2}.
\end{equation}
RMSE includes both displacement of the mean and dispersion across streams.
For scores it is computed separately in the stated parameter coordinates;
components with different units are not pooled. Likelihood entries are
\emph{log} likelihoods. An apparently close mean can conceal large errors in
individual runs, which is why both mean and RMSE are shown.
In the tables below, a proposal cell written as $u\,(v)$ means mean estimate
$u$ and RMSE $v$ relative to the displayed reference. ESS is reported as a
time-and-stream average, so it is not an accuracy error measure.

Only the LGSSM reference is exact: it is the Kalman likelihood and score.
The KSC reference is a refined grid for the implemented seven-component
observation-mixture target, not the unapproximated stochastic-volatility
model. The predator--prey reference averages two adapted Zhao--Cui local-score
fits; support, fit and differentiation-radius errors remain unresolved.
Range--bearing uses sixteen particle/Fisher reference replicates with
$65{,}536$ particles; score Monte Carlo standard errors are about $0.06502$
and $0.02268$, and finite-particle bias remains possible. The displayed SIR
reference is the rank-$40$, fit-$2$ adapted Zhao--Cui calculation. Its separate
rank-$20$, fit-$17$ check is retained when interpreting an ordering. These
references must not all be described as exact model oracles.

\IfFileExists{chapters/observation_parent_comparison_tables.tex}
 {% Generated from checked comparison.csv; do not edit numerical cells.
\begin{table}[htbp]
\centering\small
\setlength{\tabcolsep}{4pt}
\resizebox{\linewidth}{!}{%
\begin{tabular}{lrrrr}
\toprule
Model ($T$) & Reference & Original & Uniform & Weighted\\
\midrule
LGSSM (50) & -138.358 & -138.377 (0.0513889) & -138.377 (0.0523155) & -138.377 (0.0525897)\\
KSC SV (50) & -111.975 & -112.039 (0.206717) & -112.025 (0.135892) & -112.077 (0.115135)\\
Predator--prey (50) & -267.317 & -267.328 (0.145348) & -267.317 (0.148873) & -267.335 (0.17887)\\
Range--bearing (20) & 39.8555 & 39.6365 (0.331358) & 39.6231 (0.349486) & 39.7131 (0.231912)\\
SIR nominal (50) & -1669.09 & -1669.27 (0.854424) & -1669.16 (0.750419) & -1669.19 (0.654556)\\
\bottomrule
\end{tabular}}
\caption{Estimated log-likelihood mean (RMSE of the eight stream estimates)
across paired streams. Each proposal entry is reported as mean (RMSE relative
to the model-specific reference); the reference column is the reference value.
References have the model-specific limitations described in the text.}
\label{tab:parent-cross-model-likelihood}
\end{table}

\begin{table}[htbp]
\centering\small
\setlength{\tabcolsep}{4pt}
\resizebox{\linewidth}{!}{%
\begin{tabular}{lrrrr}
\toprule
Parameter & Reference & Original & Uniform & Weighted\\
\midrule
Transition & 5.86731 & 5.87984 (0.0980073) & 5.87915 (0.0993004) & 5.86609 (0.0950208)\\
Log process variance & 7.32385 & 7.32594 (0.0144067) & 7.32557 (0.0149315) & 7.32618 (0.0142106)\\
Log obs. variance & 4.13012 & 4.13771 (0.0411005) & 4.1381 (0.0414811) & 4.14015 (0.0424457)\\
Initial mean scale & 0.099617 & 0.100207 (0.00193968) & 0.100258 (0.00201672) & 0.100278 (0.00195216)\\
\bottomrule
\end{tabular}}
\caption{LGSSM: score mean (component RMSE) across eight paired streams.
Each proposal entry is mean (RMSE relative to the reference) in the displayed
parameter coordinate. Lower observed RMSE is descriptive; no cross-model score
scale is imposed.}
\label{tab:parent-cross-model-score-lgssm}
\end{table}

\begin{table}[htbp]
\centering\small
\setlength{\tabcolsep}{4pt}
\resizebox{\linewidth}{!}{%
\begin{tabular}{lrrrr}
\toprule
Parameter & Reference & Original & Uniform & Weighted\\
\midrule
Persistence & -2.0026 & -2.02203 (0.0744471) & -1.39563 (0.609442) & -1.59661 (0.407061)\\
Log obs. scale & -0.824124 & -0.698378 (0.142015) & -0.436223 (0.389578) & -0.733437 (0.0944128)\\
\bottomrule
\end{tabular}}
\caption{KSC SV: score mean (component RMSE) across eight paired streams.
Each proposal entry is mean (RMSE relative to the reference) in the displayed
parameter coordinate. Lower observed RMSE is descriptive; no cross-model score
scale is imposed.}
\label{tab:parent-cross-model-score-ksc}
\end{table}

\begin{table}[htbp]
\centering\small
\setlength{\tabcolsep}{4pt}
\resizebox{\linewidth}{!}{%
\begin{tabular}{lrrrr}
\toprule
Parameter & Reference & Original & Uniform & Weighted\\
\midrule
$r$ & -33.9911 & -34.1748 (0.302258) & -34.2382 (0.353482) & -34.3073 (0.397328)\\
Carrying capacity & -0.769203 & -0.782646 (0.0228249) & -0.783775 (0.0282161) & -0.834925 (0.0689603)\\
Half saturation & -0.00369203 & -0.00375312 (0.00139447) & -0.00405471 (0.00150476) & -0.0051175 (0.00174575)\\
$s$ & 4.04338 & 3.94216 (0.240421) & 3.98077 (0.168025) & 4.47244 (0.459779)\\
$u$ & -4.47828 & -4.4094 (0.383359) & -4.35234 (0.423594) & -4.22303 (0.424911)\\
$v$ & 4.87586 & 4.79117 (0.468782) & 4.72069 (0.51788) & 4.60354 (0.495951)\\
\bottomrule
\end{tabular}}
\caption{Predator--prey: score mean (component RMSE) across eight paired
streams. Each proposal entry is mean (RMSE relative to the reference) in the
displayed parameter coordinate. Lower observed RMSE is descriptive; no
cross-model score scale is imposed.}
\label{tab:parent-cross-model-score-predator-prey}
\end{table}

\begin{table}[htbp]
\centering\small
\setlength{\tabcolsep}{4pt}
\resizebox{\linewidth}{!}{%
\begin{tabular}{lrrrr}
\toprule
Parameter & Reference & Original & Uniform & Weighted\\
\midrule
Log process scale & -0.334787 & -0.304327 (0.634585) & -0.130327 (0.426624) & -0.252753 (0.502737)\\
Log obs. scale & -5.79515 & -5.23519 (0.758656) & -5.45566 (0.715763) & -5.394 (0.817229)\\
\bottomrule
\end{tabular}}
\caption{Range--bearing: score mean (component RMSE) across eight paired
streams. Each proposal entry is mean (RMSE relative to the reference) in the
displayed parameter coordinate. Lower observed RMSE is descriptive; no
cross-model score scale is imposed.}
\label{tab:parent-cross-model-score-range-bearing}
\end{table}

\begin{table}[htbp]
\centering\small
\setlength{\tabcolsep}{4pt}
\resizebox{\linewidth}{!}{%
\begin{tabular}{lrrrr}
\toprule
Parameter & Reference & Original & Uniform & Weighted\\
\midrule
Log infection mult. & 102.855 & 6338.21 (18648.6) & 110.512 (32.7596) & 131.782 (44.4186)\\
Log recovery mult. & -48.6029 & -2996.01 (7905.32) & -53.6253 (7.23488) & -57.2394 (13.0812)\\
Log noise mult. & -30.3076 & -5999.53 (9593.82) & -27.9792 (3.85256) & -25.5908 (4.94545)\\
\bottomrule
\end{tabular}}
\caption{SIR nominal: score mean (component RMSE) across eight paired streams.
Each proposal entry is mean (RMSE relative to the reference) in the displayed
parameter coordinate. Lower observed RMSE is descriptive; no cross-model score
scale is imposed.}
\label{tab:parent-cross-model-score-sir}
\end{table}

\begin{table}[htbp]
\centering\small
\setlength{\tabcolsep}{4pt}
\resizebox{\linewidth}{!}{%
\begin{tabular}{lrrr}
\toprule
Model & Original & Uniform & Weighted\\
\midrule
LGSSM & 957.19 & 957.143 & 957.185\\
KSC SV & 883.396 & 941.015 & 996.65\\
Predator--prey & 972.369 & 936.209 & 972.323\\
Range--bearing & 727.786 & 718.443 & 733.762\\
SIR nominal & 796.728 & 816.117 & 911.872\\
\bottomrule
\end{tabular}}
\caption{ESS averaged over time and eight streams, with $N=1008$. ESS measures weight concentration rather than likelihood or score accuracy.}
\label{tab:parent-cross-model-ess}
\end{table}
}
 {\input{../../chapters/observation_parent_comparison_tables}}

The SIR score reductions relative to the original proposal are substantial
in the measured runs. At nominal $T=50$, infection-score RMSE falls from
about $18{,}649$ to $32.76$ with uniform conditional parents and $44.42$
with weighted parents. The corresponding recovery values are $7905.32$,
$7.23$ and $13.08$; noise values are $9593.82$, $3.85$ and $4.95$.
The same dramatic statement does not transfer to every model. KSC persistence
RMSE is $0.07445$ for the original proposal, $0.60944$ for uniform parents
and $0.40706$ for weighted parents. In predator--prey, weighted parents have
larger observed RMSE than the original in all six score components. These
are adverse accuracy results even though the runs pass their recorded
numerical-validity checks.

Between the two conditional choices, uniform parents have lower observed
score RMSE in both range--bearing components, five of six predator--prey
components and all three nominal SIR components. Weighted parents have lower
observed score RMSE in both KSC components and three of four LGSSM
components. Weighted parents have lower log-likelihood RMSE in KSC,
range--bearing and nominal SIR; uniform parents have lower values in LGSSM
and predator--prey. These mixed comparisons make uniform parents a serious
alternative when score accuracy is the priority.

Across the twenty SIR parameter--horizon cells with usable references,
Table~\ref{tab:parent-sir-allocation-counts} requires the same RMSE ordering
against both reference fits. A reference-dependent ordering is left unresolved.
At nominal $T=50$, the first and second reference score vectors are
\begin{align*}
  s^{(1)} &= (102.85487,-48.60288,-30.30761),\\
  s^{(2)} &= (108.57912,-50.42205,-30.38459).
\end{align*}
Uniform parents have lower observed score RMSE in all three components
against either fit.

\begin{table}[htbp]
\centering\small
\begin{tabular}{lrrr}
\toprule
Quantity & Uniform lower & Weighted lower & Reference dependent\\
\midrule
Log likelihood & 5 & 15 & 0\\
Infection score & 13 & 3 & 4\\
Recovery score & 15 & 2 & 3\\
Noise score & 13 & 7 & 0\\
\bottomrule
\end{tabular}
\caption{Descriptive RMSE counts over twenty usable SIR cells. The counts do
not constitute a statistical ranking or a probability of superiority.}
\label{tab:parent-sir-allocation-counts}
\end{table}

Parent weighting has a more consistent benefit for weight balance in the SIR
campaign. In every usable cell, its time-mean and time-minimum ESS exceed
those of uniform conditional parents in all eight paired streams: $160/160$
pairs for each summary. The increase in stream-averaged time-mean ESS ranges
from $6.7\%$ to $102.7\%$. The direction also holds in all $33$ executed
cells, or $264/264$ paired streams, including cells whose reference fails the
accuracy screen. ESS can be compared there without relying on that reference.
The KSC and range--bearing extensions also show higher time-mean ESS for
weighted parents in all eight paired streams. This does not establish higher
ESS at every time step or on every future dataset. In particular, the
cross-model mean ESS values do not make either conditional proposal uniformly
better than the original.

The retained default is therefore an engineering choice for improved weight
balance: observation-conditioned proposals with observation-weighted parents.
Uniform conditional parents remain an explicit, competitive alternative.
The default is not an accuracy-optimality claim. Both choices are markedly
more accurate than the original in the nominal SIR score experiment; the
KSC and predator--prey counterexamples prevent extending that statement to
all models and parameters.

All reported arms remain numerically viable under their saved checks.
The five-model comparison establishes no statistical ranking. Eight streams
on one dataset, approximate references and inherited controls limit its
interpretation. The SIR analysis retains two exploratory componentwise
findings that favor opposite parent rules, without a multiplicity correction;
its large original-versus-conditional score differences do not pass the full
paired mean-squared-error criterion. The next evidence for an accuracy choice
is stronger references where needed, scope-specific calibration on disjoint
data, and predeclared paired uncertainty across independent datasets.

The numerical source is the checked CSV and source-hash manifest under
\path{docs/plans/artifacts/observation-parent-all-model-comparison-20261028/}.
The accompanying result note retains the target definitions, per-stream
checks, inference limits and links to each completed campaign. The tables
here are generated directly from that CSV so that the research note and
monograph report the same measurements.


\section{Why the failures matter for this inference problem}
The failures concern different mathematical quantities. A proposal can put
particles in the wrong place, an importance correction can leave highly
variable weights, and an accurately evaluated derivative can still estimate
the model score poorly. Treating these as one problem obscures which repair
an experiment has actually tested.

\subsection{A correct flow applied to the wrong conditional distribution}

A UKF describes predictive uncertainty using moments. An LEDH map then uses
the observation to move particles. Neither fact guarantees that the map's
reference covariance matches the distribution of the particles passed to it.
The distinction is particularly consequential when the predictive spread
comes mainly from differences between ancestors rather than process noise
within one ancestor.

A scalar Gaussian calculation shows the issue without numerical flow error.
Suppose a map is constructed for prior $\mathcal N(m,P)$ and observation
$y=x+\varepsilon$, $\varepsilon\sim\mathcal N(0,R)$. An exact affine
Gaussian transport is
\begin{equation}
 F_{m,P}(x)=m+\frac{P}{P+R}(y-m)
              +\sqrt{\frac{R}{P+R}}(x-m).
\end{equation}
If its input really has law $\mathcal N(m,P)$, the output has the correct
posterior mean and variance. If the input is instead a particular parent's
transition $X\sim\mathcal N(a,Q)$, direct substitution gives
\begin{align}
 \mathbb E[F_{m,P}(X)]&=m+\frac{P}{P+R}(y-m)
                    +\sqrt{\frac{R}{P+R}}(a-m),\\
 \operatorname{Var}[F_{m,P}(X)]&=\frac{RQ}{P+R}.
\end{align}
The desired conditional posterior has mean
$a+Q(y-a)/(Q+R)$ and variance $QR/(Q+R)$. For $a=m=0$, $Q=1$,
$P=100$, $R=100$ and $y=10$, the map yields mean $5$ and variance $0.5$,
whereas conditioning the actual transition yields approximately $0.0990$
and $0.9901$. The map is correct for its stated Gaussian input; using it as
the exact conditional update for this different input is wrong.

This example identifies a mechanism, rather than measuring how much of a
particular nonlinear run's error it explains. A proposal with this mismatch
can still have a valid importance correction if its actual density is
evaluated. But valid weights cannot guarantee that a finite sample covers
the important region well. Parent-conditioned proposals address the
within-parent mismatch; predictive-observation weighting addresses the
allocation of children between parents. These are distinct repairs.

\subsection{An exact update of an inaccurate incoming mixture}

Even exact conditioning only updates the incoming particle approximation.
If the previous cloud omits a region of the filtering distribution, the
finite-mixture normalizer $N^{-1}\sum_j L_j(y)$ can differ from the true
predictive likelihood. Exact conditional sampling cannot reconstruct
ancestors that are absent. Transport and moment restoration also retain only
specified properties of the cloud; preserving its first two moments does
not determine its tails, multimodality or parameter sensitivity.

The familiar weight diagnostic
\begin{equation}
 \operatorname{ESS}=\frac{1}{\sum_i \bar w_i^2}
\end{equation}
measures concentration of normalized importance weights. With exact
observation-weighted conditional sampling of the finite mixture, the
incremental weights are constant and ESS is $N$. This does not make the
incoming approximation exact, nor does it imply $N$ independent informative
ancestors after earlier resets. Higher ESS has practical value because less
of the represented mass is carried by a few particles. It remains a measure
of weight balance, not of distance to the model posterior or its score.

The SIR controls make this limitation observable. In the pure
conditional-only ablation (the defensive branches are removed), at nominal
$T=50$ the weighted proposal attains ESS $1008$, but its infection-score
RMSE is $95.82$--$100.96$ across the two references. The conditional mixture
with uniform parents gives $31.91$--$32.76$ in the corresponding
three-branch comparison. At noise multiplier $1.25$,
$T=50$, identity-only and conditional-only uniform proposals also have lower
observed noise-score RMSE than the weighted conditional mixture. These
controls prevent a claim of accuracy optimality. They do not remove the
value of the measured reduction in weight concentration.

\subsection{Differentiating the finite calculation is a different test}

The model score is $s(\theta)=\nabla_\theta\log p_\theta(y_{1:T})$.
The implemented analytical calculation differentiates the finite filter
value $\widehat\ell_N(\theta;u)$ for fixed random design $u$, with sampled
categorical labels held locally fixed. A finite-difference comparison away
from selection boundaries checks that branchwise derivative. It does not
test whether the derivative is close to $s(\theta)$, nor whether it equals
the derivative after averaging over parameter-dependent selections.

A two-parent example makes the distinction explicit. Let $U$ be uniform on
$(0,1)$, choose parent 1 if $U<\alpha(\theta)$ and parent 2 otherwise, and
return fixed values $a$ and $b$. For almost every fixed $U$, the selected
value is locally constant in $\theta$, so its branchwise derivative is zero.
Yet
\begin{equation}
 \mathbb E[h_\theta(U)]=\alpha(\theta)a+
               [1-\alpha(\theta)]b,
 \qquad
 \frac{\mathrm d}{\mathrm d\theta}\mathbb E[h_\theta(U)]
       =\alpha'(\theta)(a-b).
\end{equation}
The derivative of the expectation contains the change in selection
probability. This example demonstrates why local derivative parity alone
cannot settle a model-score claim. It does not prove that every observed
error in the full filter is caused by a missing selection contribution:
the full calculation also differentiates densities, normalizations, state
locations and resets, and some terms can cancel under particular identities.
That causal claim requires an isolated comparison.

There is a second distinction even without categorical choices. If a positive
likelihood estimator satisfies $\mathbb E[\widehat Z]=Z$, and differentiation
may pass through the expectation, then
$\mathbb E[\nabla\widehat Z]=\nabla Z$. It does not follow that
\begin{equation}
 \mathbb E\!\left[\frac{\nabla\widehat Z}{\widehat Z}\right]
       =\frac{\nabla Z}{Z}.
\end{equation}
The random denominator is correlated with the numerator. Correctness of an
unnormalized derivative identity therefore does not remove finite-particle
score bias. Nor does a small error in a likelihood value bound its derivative
error. These are reasons to compare likelihood and score separately.

\subsection{What the earlier score repairs did and did not solve}

The observation-sidecar and kernel-mixture investigation exposed a target
distinction before parent conditioning was introduced. An expected
conditional log likelihood is generally different from the log marginal
likelihood. For a fixed mixing measure, Jensen's inequality gives
$\mathbb E[\log g_\theta(y\mid X)]\leq
\log\mathbb E[g_\theta(y\mid X)]$; their derivatives need not agree.
An unbiased gradient identity for the former cannot establish an unbiased
model score for the latter. The research note derives the precise
importance-weighted gradient identity and its normalization requirements.

The subsequent complete kernelized-observation recursion and
all-components resampling reference corrected more of the finite
calculation and passed their declared derivative checks. Nevertheless, the
calibrated resampling construction had higher observed model-score MSE than
the canonical comparator in both tested matrix-Gaussian scopes. That is a
negative result for the tested construction, rather than evidence that the
derivative check failed or that every kernel-mixture proposal must fail.
The later simulation-supervised and neural-guide comparisons at $N=64$
were withdrawn; they cannot support an accuracy ranking or rejection of
guidance. The full derivations and scope-specific evidence are retained in
the research note, rather than treating those discarded runs as support
for the present default.

\subsection{The model-specific failures still requiring explanation}

KSC is a useful counterexample to an ESS-based score argument. Weighted
parents raise mean ESS from about $883$ to $997$ and reduce log-likelihood
RMSE, yet persistence-score RMSE increases from about $0.074$ to $0.407$.
Because the conditional calculation is exact for the implemented
seven-mixture observation target, a generic claim that all remaining error
comes from an unconditioned child is inadequate. Recursive representation,
selection sensitivity and inherited numerical controls remain possible
causes; the completed experiment does not isolate their contributions.

In range--bearing, conditioning is locally Gaussian and the bearing residual
is wrapped. Its conditional proposal is therefore an approximation to a
nonlinear observation update. Weighted parents reduce the observed
process-scale score RMSE relative to the original, but increase the
observation-scale score RMSE; uniform parents have lower RMSE in both
components than weighted parents. Reference Monte Carlo error limits the
precision of this comparison. The experiment does not establish whether
local linearization, allocation, or recursive sensitivity explains the
remaining discrepancy.

Predator--prey is the clearest warning against a universal cross-model
claim: all six weighted-parent score RMSEs exceed the original comparator's
values. The adapted Zhao--Cui local fits are themselves approximate, so
better reference and support checks are needed before assigning the full
gap to the particle proposal. LGSSM supplies the exact-oracle control:
the small, mixed differences show no broad accuracy advantage for either
parent rule. SIR supplies the opposite extreme, where conditioning the
children removes very large observed score errors but additional parent
weighting often trades score accuracy for ESS and likelihood accuracy.

The next discriminating experiment should hold the incoming cloud fixed,
compare exact finite-mixture integration with conditional sampling, and
separate local derivatives from derivatives of expectations over selection.
Only after that distinction is resolved should independently calibrated
recursive runs identify the effect of reset and allocation choices. The
current evidence supports retaining weighted parents as the requested
default and exposing uniform parents as a competitive option. It does not
close model-score correctness, HMC readiness, or accuracy-optimality questions.


\section{Scope of the parent-selection default}
The comparison above concerns the Gaussian conditional-proposal route used in
the reported model experiments. Deterministic structural-completion models
use a separate innovation-space proposal, whose implementation currently
selects parents uniformly. The observation-weighted rule is not implemented
for that route: its parent probabilities would need to be computed from the
innovation-space predictive observation law, with the resulting marginal
proposal density retained in the importance correction. The code therefore
rejects the weighted-parent setting for structural completion rather than
silently running the uniform rule under a weighted-policy label.
